%\documentclass[11pt,a4paper,uplatex,dvipdfmx]{ujarticle} 		% for uplatex
\documentclass[11pt,a4j,dvipdfmx]{jarticle} 					% for platex
\input{pieces/form00_header} % pieces
\input{pieces/kakenhi7} % pieces
\input{pieces/form01_header} % pieces
\input{pieces/form02_2023_header} % pieces
\input{pieces/hook3} % pieces
%#Name: gakuhen_a_keikaku
\input{pieces/form04_jsps_headers} % pieces
\input{pieces/form04_mext_headers_2020} % pieces
\input{pieces/form04_gakuhen_a_keikaku_header} % pieces
% ===== Global definitions for the Kakenhi form ======================
%　基本情報
%
%------ 研究課題名  -------------------------------------------
\newcommand{\研究課題名}{象の卵}

%----- 研究機関名と研究代表者の氏名-----------------------
\newcommand{\研究機関名}{逢坂大学}
\newcommand{\研究代表者氏名}{湯川秀樹}
\newcommand{\me}{\underline{\underline{H.~Yukawa}}} 
%---- 研究期間の最終年度 ----------------
\newcommand{\研究期間の最終元号年度}{9}  %令和で、半角数字のみ
%========================================

\input{pieces/inst_general} % pieces
\input{pieces/inst_GHa_keikaku} % pieces
\input{keikaku_defs}
% user07_header
% ===== my favorite packages ====================================
% ここに、自分の使いたいパッケージを宣言して下さい。
\usepackage{wrapfig}
%\usepackage{amssymb}
%\usepackage{mb}
%\DeclareGraphicsRule{.tif}{png}{.png}{`convert #1 `dirname #1`/`basename #1 .tif`.png}
\usepackage{lineno}

% ===== my personal definitions ==================================
% ここに、自分のよく使う記号などを定義して下さい。
\newcommand{\klpionn}{K_L \to \pi^0 \nu \overline{\nu}}
\newcommand{\kppipnn}{K^+ \to \pi^+ \nu \overline{\nu}}

% ----- 業績リスト用 -------------
\newcommand{\paper}[6]{%
	% paper{title}{authors}{journal}{vol}{pages}{year}
	\item ``#1'', #2, #3 {\bf #4}, #5 (#6).			% お好みに合わせて変えてください。
}

\newcommand{\etal}{\textit{et al.\ }}
\newcommand{\ca}[1]{*#1}	% corresponding author;   \ca{\yukawa}  みたいにして使う
\newcommand{\invitedtalk}{招待講演}

\newcommand{\yukawa}{H.~Yukawa}					% no underline
%\newcommand{\yukawa}{\underline{\underline{H.~Yukawa}}}	% with 2 underlines
\newcommand{\tomonaga}{S.~Tomonaga}

\newcommand{\prl}{Phys.\ Rev.\ Lett.\ }		% よく使う雑誌も定義すると楽

% ===== 欄外メモ ==================
\newcommand{\memo}[1]{\marginpar{#1}}
%\renewcommand{\memo}[1]{}	% 全てのメモを表示させないようにするには、行頭の"%"を消す

%%\input{../../pieces/user07_header}	% expand
%\input{../../sample/simple/contents}	% skip
\input{pieces/hook5} % pieces

\begin{document}
\input{pieces/hook7} % pieces
%#Split: 01_purpose_plan  
%#PieceName: p01_purpose_plan
\input{pieces/p01_purpose_plan_00}
\section{ （ａ）研究計画、研究方法など}
%　　　　＜＜最大　6ページ＞＞

%s02_purpose_plan_with_abstract_gakuhen
\noindent
\textbf{（概要）}\\
%begin 研究目的及び研究計画の概要空行付き ＝＝＝＝＝＝＝＝＝＝＝＝＝＝＝＝＝＝＝＝
	象の卵の研究の目的は．．．	

	象の卵の研究計画と方法は．．．
	\vspace*{10zw}	% （概要）と（本文）の間が10行程度になるよう、必要に応じて値を調整してください。	
%end 研究目的及び研究計画の概要空行付き ＝＝＝＝＝＝＝＝＝＝＝＝＝＝＝＝＝＝＝＝

\noindent
\rule{\linewidth}{1pt}\\
\noindent
\textbf{（本文）}
%begin 研究目的と研究計画	＝＝＝＝＝＝＝＝＝＝＝＝＝＝＝＝＝＝＝＝
\JSPSInstructions	% <-- 留意事項。これは消すか、コメントアウトしてください。

\textbf{\\　　　　　＊＊＊　以下は、あくまで例です。真似しないでください。　＊＊＊\\
　　　　　＊＊＊　本文はもちろん、節の切り方や論理の組み方は　　　＊＊＊\\
　　　　　＊＊＊　ご自分の気に入ったスタイルで書いてください。　　＊＊＊}

\subsection{計画研究の名前}
        領域全体の計画書を書くときに用いた\texttt{keikaku\_defs.tex}で定義したコマンド
        (例えば
        \texttt{\textbackslash codeZoo, \textbackslash titleZoo, \textbackslash codeNicknameZoo}など)を用い、
        各計画研究の記号や研究課題名やその略称
        （\codeZoo, \titleZoo, \codeNicknameZoo)
        を表示するのが楽です。
       書き間違えないし、あとで名前が変わっても簡単に直せます。

	 象の卵の研究目的は．．．

	象の卵の研究計画は．．．

	\input{blahblah}  % << only for demonstration. Please delete it or comment it out.
	\input{blahblah}  % << only for demonstration. Please delete it or comment it out.
		

	\vspace{1cm}
	\begin{thebibliography}{99}
		\bibitem{teramura} 寺村輝夫、「ぼくは王様 - ぞうのたまごのたまごやき」.
	\end{thebibliography}
%end 研究目的と研究計画	＝＝＝＝＝＝＝＝＝＝＝＝＝＝＝＝＝＝＝＝

\input{pieces/p01_purpose_plan_01}

%#Split: 07_rights  
%#PieceName: p07_rights
\input{pieces/p07_rights_00}
\section{人権の保護及び法令等の遵守への対応}
%　　　　＜＜最大　1ページ＞＞

% s09_rights
%begin 人権の保護及び法令等の遵守への対応 ＝＝＝＝＝＝＝＝＝＝＝＝＝＝＝＝＝＝＝＝
	象の卵のES細胞の培養、象のクローンの生成などは行わない。

	\LaTeX の便利な機能については、\texttt{egg\_***.tex} や\texttt{sample\_pdf/egg\_***.pdf}をご覧ください。
%end 人権の保護及び法令等の遵守への対応 ＝＝＝＝＝＝＝＝＝＝＝＝＝＝＝＝＝＝＝＝

\input{pieces/p07_rights_01}

%#Split: members_info % keep

% 研究代表者の調書 (not 長所)　＝＝＝＝＝＝＝＝＝＝＝＝＝＝＝
\KLInput{daihyo_info}

% 研究分担者の調書　＝＝＝＝＝＝＝＝＝＝＝＝＝＝＝
\KLInput{buntansha_info}

% 研究分担者を追加する場合は、buntansha_info.tex を複製して別の名前にし、
% この下に追加してください。
% \input ではなく、\KLInput を用いると、読み込んだファイルの中で定義した
% マクロはそのファイルの中でのみ有効となり、他のファイルでの定義と干渉しません。
%\KLInput{taro_info}
%\KLInput{hanako_info}
%\KLInput{jiro_info}
%\KLInput{saburo_info}


%#Split: 99_tail
\input{pieces/hook9} % pieces
\end{document}

